\documentclass[slides]{pgnotes}

\title{Regular expressions (Regex)}

\begin{document}

\maketitle

\section{Regular expressions (regex)}\label{regex}

Tool to search text for specific patterns:
\begin{itemize}
\item Very flexible.
\item Similar implementations in many languages, slightly different dialects.
\item Python has the \texttt{re} package as standard.
\item Supported by PostgreSQL using the \texttt{\textasciitilde{}} operator.
\end{itemize}
  
\section{Building a regex}\label{building-a-regex}

\subsection{Anchors}\label{anchors}

Without an anchor, \texttt{\^{}} or \texttt{\$}, the regex will return true if the specified text exists.

\begin{itemize}
\item Use a preceeding \texttt{\^{}} for matches \emph{starting with} the specified pattern.
\item Use a terminating \texttt{\$} for matches \emph{ending with} the specified pattern.
\item Use both \texttt{\^{}} and \texttt{\$} to match a specific string exactly. (Could just use \texttt{LIKE} or \texttt{=}!)
\end{itemize}

\section{PostgreSQL}\label{postgresql}

\subsection{Regex matching operators}\label{regex-matching-operators}

The Regex matching operators allow regex matching in queries:

\begin{verbatim}
~    regex matches, case sensitive
~*   regex matches, case insensitive
!~   regex does not match, case sensitive
!~*  regex does not match, case insensitive
\end{verbatim}

\subsection{CHECK constraint}\label{check-constraint}

The CHECK constraint can force a TEXT column to conform to a particular regex:

\begin{verbatim}
CHECK (columnname ~ 'regex_expression_here')
\end{verbatim}

\subsection{DOMAIN}\label{domain}

A DOMAIN allows us to define a particular datatype in terms of:

\begin{enumerate}
\def\labelenumi{\arabic{enumi}.}
\tightlist
\item
  A base type
\item
  A CHECK constraint
\end{enumerate}

Pattern:

\begin{verbatim}
CREATE DOMAIN domainname AS TEXT CHECK ( VALUE ~ 'regex_here' )
\end{verbatim}

Generally NOT a good idea to force NOT NULL on a DOMAIN as limits
flexibility.

\end{document}

